% Wrapper pour changer le nombre de colonnes d'un chant individuel.
\newcommand{\songcols}[1]{\songcolumns{#1}}

% Page de séparation thématique — papier déchiré + pictos
% #1 = titre de la section
% #2 = image du papier (ex: images/papier-vert.png)
% #3 = bloc tikz de pictos (nœuds tikz pré-générés)
% Décalage du contenu pour centrer dans la zone de texte (marges asymétriques)
\newlength{\sectionpageshift}
\setlength{\sectionpageshift}{0.5\dimexpr\oddsidemargin-\evensidemargin\relax}
\newcommand{\songsectionpage}[3]{%
  \cleardoublepage
  \thispagestyle{empty}%
  \markboth{#1}{#1}%
  \begin{tikzpicture}[remember picture, overlay]
    % Centre de la zone de texte (décalé par rapport au centre de la page)
    \coordinate (contentcenter) at ([xshift=\sectionpageshift]current page.center);
    % Fond de page section (couvre toute la page)
    \fill[fill={rgb,255:red,243;green,236;blue,189}] (current page.south west) rectangle (current page.north east);
    % Papier déchiré en bandeau horizontal (image portrait tournée 90°)
    \node[anchor=center, rotate=90, inner sep=0pt]
      at (contentcenter)
      {\includegraphics[width=0.65\paperheight]{#2}};
    % Pictos par-dessus le bandeau
    #3
    % Titre en blanc par-dessus tout
    \node[anchor=center,
          font=\fontsize{36}{42}\selectfont\bfseries\rmfamily\scshape,
          text=white]
      at (contentcenter) {#1};
  \end{tikzpicture}%
  \clearpage
}

% Variante pour image de bandeau déjà en paysage (pas de rotation)
\newcommand{\songsectionpageh}[3]{%
  \cleardoublepage
  \thispagestyle{empty}%
  \markboth{#1}{#1}%
  \begin{tikzpicture}[remember picture, overlay]
    \coordinate (contentcenter) at ([xshift=\sectionpageshift]current page.center);
    \fill[fill={rgb,255:red,243;green,236;blue,189}] (current page.south west) rectangle (current page.north east);
    % Papier déchiré — image déjà en paysage, pas de rotation
    \node[anchor=center, inner sep=0pt]
      at (contentcenter)
      {\includegraphics[width=0.95\paperwidth]{#2}};
    #3
    \node[anchor=center,
          font=\fontsize{36}{42}\selectfont\bfseries\rmfamily\scshape,
          text=white]
      at (contentcenter) {#1};
  \end{tikzpicture}%
  \clearpage
}

% --- Pastilles colorées de section dans l'index ---
% Couleurs extraites des bandeaux de papier déchiré
\definecolor{idxsec1}{HTML}{4EC983}  % SGDF (papier-vert)
\definecolor{idxsec2}{HTML}{EE8800}  % Bénédicités (papier-orange)
\definecolor{idxsec3}{HTML}{2196F3}  % Coin du feu (papier-bleu)
\definecolor{idxsec4}{HTML}{D32F2F}  % Chansons du monde (papier-rouge)
\definecolor{idxsec5}{HTML}{00695C}  % Temps spi (papier-vert compa)

\makeatletter
% Extraction du numéro de section par pattern-matching direct
% \songlink{songN-...}{page} → on extrait N
\def\idx@grabsec\songlink#1#2#3\idx@stop{\idx@colorfromsec#1\idx@stop}
\def\idx@colorfromsec song#1-#2\idx@stop{%
  \mbox{\raisebox{0.2ex}{\tikz{\node[circle, fill=idxsec#1, minimum size=4.5pt, inner sep=0pt] {};}}}%
  \hspace{1pt}%
}

% Dessiner la pastille colorée avant le titre
\newcommand{\idxsecdisc}[1]{%
  \idx@grabsec#1\idx@stop
}

% Redéfinir le rendu de l'index des titres pour ajouter les pastilles
\renewcommand\SB@maketitleindex{%
  \ifnum\idxheadwidth>\z@%
    \renewenvironment{SB@lgidx}[1]{
      \hbox{\SB@colorbox\idxbgcolor{\vbox{%
        \hbox to\idxheadwidth{{\idxheadfont\relax##1}\hfil}%
      }}}%
      \nobreak\vskip3\p@\@plus2\p@\@minus2\p@\nointerlineskip%
    }{\penalty-50\vskip5\p@\@plus5\p@\@minus4\p@}%
  \else%
    \renewenvironment{SB@lgidx}[1]{}{}%
  \fi%
  \renewenvironment{SB@smidx}[1]{}{}%
  \renewcommand\idxentry[2]{%
    \SB@ellipspread{\idxsecdisc{##2}\idxtitlefont\relax\ignorespaces##1\unskip}%
                   {{\idxrefsfont\relax##2}}%
  }%
  \renewcommand\idxaltentry[2]{%
    \SB@ellipspread{\idxlyricfont\relax\ignorespaces##1\unskip}%
                   {{\idxrefsfont\relax##2}}%
  }%
  \SB@displayindex%
}
\makeatother
